%%%
%
% $Autor: Wings $
% $Datum: 2021-05-14 $
% $Pfad: ArduinoSerial $
% $Dateiname: 
% $Version: 4620 $
%
% !TeX spellcheck = en_GB/de_D
% !TeX program = pdflatex
% !BIB program = biber/bibtex
% !TeX encoding = utf8
%
%%%



\chapter{Serielle Kommunikation zwischen PC und Arduino} \label{SerialCom}

Um eine Kommunikation zwischen dem PC und dem Arduino Nano 33 BLE Sense über USB mit Python zu realisieren, bietet sich die Verwendung des seriellen Protokolls an.

In diesem Kapitel wird anhand eines einfachen Beispiels gezeigt, wie über die serielle Schnittstelle eine RGB-LED gesteuert werden kann.

Die Umsetzung erfolgt in fünf Schritten:

\begin{enumerate}
	\item Installation einer Bibliothek auf dem PC
	\item Entwicklung des Python-Programms auf dem PC
	\item Vorbereitung des Arduino-Boards
	\item Entwicklung des Arduino-Sketches
	\item Einheitliche Konfiguration der seriellen Schnittstelle
\end{enumerate}

\section{Installation einer Bibliothek auf dem PC}

Das serielle Protokoll ermöglicht den Datenaustausch zwischen Geräten über eine serielle Schnittstelle. Unter Windows kann zur Kommunikation mit seriellen Geräten das Python-Paket \PYTHON{PySerial} verwendet werden.

Im ersten Schritt ist das Paket \PYTHON{PySerial} mithilfe des Programms \SHELL{pip} zu installieren:

\medskip

\SHELL{pip install pyserial}

\section{Entwicklung des Python-Programms auf dem PC}

Nach der Installation des Pakets \FILE{PySerial} kann die Klasse \PYTHON{serial} verwendet werden. Die Methode \PYTHON{serial.Serial} erzeugt ein Objekt, das die Verbindung zum Arduino repräsentiert. Dabei müssen der Portname (z.B. COM7) und die Baudrate (z.B. 9600) mit denen im Arduino-Sketch übereinstimmen. Über die Methoden \PYTHON{write} und \PYTHON{read} können Daten zum Arduino gesendet bzw. von diesem empfangen werden. Das folgende Beispiel demonstriert das Senden eines Zeichens an den Arduino und das Empfangen einer Antwort:

{
	\captionof{code}{Python-Programm für den Windows-PC zur Kommunikation mit einem Arduino über die serielle Schnittstelle.}
	\label{Nano:TestSerialPC}
	\PythonExternalO{../../Code/Nano33BLESense/Serial/TestSerialPC.py}
}

\section{Entwicklung des Arduino-Sketches und Vorbereitung des Arduino}

{
	\captionof{code}{Arduino-Sketch zur Kommunikation mit einem Windows-PC über die serielle Schnittstelle.}
	\label{Nano:TestSerial}
	\ArduinoExternal{}{../../Code/Nano33BLESense/Serial/TestSerial.ino}
}

\bigskip

\begin{itemize}
	\item Das Arduino-Board wird über ein USB-Kabel mit dem Computer verbunden.
	\item Der gewünschte Sketch wird auf das Arduino geladen.
\end{itemize}

\section{Einheitliche Konfiguration der seriellen Schnittstelle}

Folgende Anpassungen sind erforderlich:

\begin{itemize}
	\item Der Port \PYTHON{ArduinoPort} ist entsprechend der Arduino-Verbindung anzupassen.
	\item Die Baudrate \PYTHON{Baudrate} muss mit der im Arduino-Sketch verwendeten Baudrate übereinstimmen.
\end{itemize}

\section{Weiterführende Literatur}

\begin{itemize}
	\item
	\HREF{https://www.instructables.com/Python-Serial-Port-Communication-Between-PC-and-Ar/}{Python Serial Port Communication Between PC and Arduino Using PySerial Library}
	
	pgsql
	Kopieren
	\item \HREF{https://arduino.stackexchange.com/questions/4891/how-would-i-establish-a-serial-connection-with-python-without-using-the-serial-m}{How would I establish a serial connection with Python without using the serial monitor?}
\end{itemize}

\bigskip

\textbf{Verwendung des seriellen Monitors:}

\begin{itemize}
	\item Der serielle Monitor der Arduino-IDE kann geöffnet werden, wobei darauf zu achten ist, dass dort dieselbe Baudrate eingestellt ist.
	\item Über den seriellen Monitor lassen sich Daten zwischen dem Arduino und dem Computer senden und empfangen.
\end{itemize}

\bigskip

Ein Beispiel für ein Python-Programm, mit dem sich die RGB-LED des Arduino Nano 33 BLE Sense mithilfe des Pakets \PYTHON{pySerial} steuern lässt, ist \HREF{https://arduino.stackexchange.com/questions/91319/serial-communication-between-python-and-arduino-nano-ble-sense-33-for-running-si}{hier} zu finden. Der zugehörige Arduino-Sketch, der die serielle Bibliothek verwendet, befindet sich unter folgendem Link: \HREF{https://micropython.org/download/ARDUINO_NANO_33_BLE_SENSE/}{hier}.

\section{Beispiel: Übertragung eines Bildes über die serielle Schnittstelle}

Dieses Beispiel ist in der zweiten Auflage des TinyML Cookbooks dokumentiert. \cite{Iodice:2023}, ab Seite 361.

